\chapter{Custom Gradients for LEDH-PFPF-OT}
\label{ch:bf-ledh-pfpf-ot-custom-gradient}

The preceding chapters introduced three ingredients separately: local exact
Daum--Huang particle flow, proposal-corrected PF-PF weighting, and
entropy-regularized OT resampling.  In the OT sequence immediately before this
chapter, the deterministic baseline fixed the transport-object and barycentric-output
contract, while the entropic chapter fixed the exact teacher object, the finite
Sinkhorn route that approximates it, and the rule that later differentiation claims
must attach to that executed finite route.  This chapter explains how those pieces can
be differentiated without asking an automatic-differentiation system to retain
the entire filtering graph in memory.  The key object is a custom
vector-Jacobian product, or VJP, for the executed LEDH-PFPF-OT scalar.

The chapter is intentionally narrow.  It does not claim that the relaxed
OT-resampled particle filter is the same object as categorical resampling.  It
does not claim that finite Sinkhorn iteration is the exact unregularized OT
optimizer.  It does not claim HMC correctness for a nonlinear model.  Its claim
is more basic and more useful: if the value path is a declared finite
LEDH-PFPF-OT computation, then the gradient path must be the derivative of that
same computation, and the expensive OT block has a structured adjoint.

\section{The scalar whose derivative is wanted}
\label{sec:bf-ledh-ot-scalar}

Let $\theta$ denote the model parameter and let all random numbers used by the
diagnostic or sampler target be fixed as part of the deterministic value path.
For one filtering pass, write the executed scalar as
\begin{equation}
\label{eq:bf-ledh-ot-executed-scalar}
  \mathcal L_K(\theta)
  =
  \operatorname{Reduce}
  \left(
    \operatorname{LEDH\mbox{-}PFPF\mbox{-}OT}_K
    (\theta; y_{1:T}, \xi)
  \right),
\end{equation}
where $\xi$ denotes the fixed particle and process-noise streams, $K$ denotes
the finite Sinkhorn or annealed-transport iteration budget, and
\(\operatorname{Reduce}\) is the declared scalar reduction, such as a mean over
fixed seeds.  The derivative needed by HMC or optimization is
\[
  \nabla_\theta \mathcal L_K(\theta),
\]
not the derivative of a different resampling rule, a different Sinkhorn stopping
policy, or a different random seed construction.

\begin{definition}[Same-scalar custom-gradient contract]
\label{def:bf-ledh-ot-same-scalar}
A custom-gradient implementation for LEDH-PFPF-OT is same-scalar if its primal
return is exactly \(\mathcal L_K(\theta)\) and its VJP returns
\[
  v^\top D_\theta \mathcal L_K(\theta)
\]
for every upstream scalar multiplier \(v\).  For a scalar log target the usual
case is \(v=1\).
\end{definition}

This contract is the same HMC principle as Chapter~\ref{ch:bf-custom-gradient-wrappers},
but the present chapter specializes it to the LEDH-PFPF-OT computation.  The
particle-flow part follows the Li--Coates invertible-particle-flow proposal
logic \citep{li2017particle}; the OT relaxation follows the entropy-regularized
resampling route of \citet{corenflos2021differentiable}, with the Sinkhorn
background of \citet{cuturi2013sinkhorn} and
\citet{peyre2019computational}.
The derivative identities below are project derivations for the declared
BayesFilter finite program; the citations explain the particle-flow,
entropy-regularized OT, and Sinkhorn objects being differentiated.

\section{Forward decomposition}
\label{sec:bf-ledh-ot-forward-decomposition}

At one observation time, the value path can be written as a composition
\begin{equation}
\label{eq:bf-ledh-ot-forward-composition}
  \theta
  \longmapsto
  (X^-, \log w^-, A_{\mathrm{flow}})
  \longmapsto
  (X, \log w)
  \longmapsto
  \widetilde X
  \longmapsto
  \ell_t .
\end{equation}
Here \(X^-\in\mathbb R^{N\times d}\) is the pre-observation proposal cloud,
\(\log w^-\) are proposal-correction log weights, and
\(A_{\mathrm{flow}}\) denotes the LEDH auxiliary quantities needed for the
proposal density and log-determinant correction.  The local LEDH construction
uses particle-local linearizations and particle-local covariance state as in
Li--Coates PF-PF.  The OT block then maps the weighted cloud to an equal-weight
cloud
\[
  (X,\log w) \longmapsto \widetilde X.
\]

This chapter focuses on the last map because it is the part that can dominate
memory.  Naive reverse-mode autodiff through all pairwise costs, Sinkhorn
updates, and barycentric products may retain objects of size \(O(KN^2)\).  A
custom VJP instead treats the OT block as a finite differentiable program and
applies its adjoint by recomputation, streaming reductions, or checkpointing.

\section{Transport orientation}
\label{sec:bf-ledh-ot-orientation}

The OT chapters used a coupling \(\Pi\) with source index first and output slot
second.  The TensorFlow LEDH-PFPF-OT route is more naturally described by the
row-stochastic transport matrix
\begin{equation}
\label{eq:bf-ledh-ot-row-transport}
  T = N\Pi^\top,
  \qquad
  T\mathbf 1=\mathbf 1,
  \qquad
  T^\top\mathbf 1=Nw,
\end{equation}
where \(w=\operatorname{softmax}(\log w)\).  The transported equal-weight cloud
is
\begin{equation}
\label{eq:bf-ledh-ot-barycentric-output}
  \widetilde X = T X,
  \qquad
  \widetilde x_j=\sum_{i=1}^N T_{ji}x_i .
\end{equation}
Thus row \(j\) of \(T\) tells how output slot \(j\) averages the source
particles.  The row sums make each output particle a barycenter.  The column
sums say that source particle \(i\) contributes total mass proportional to its
weight.

\begin{remark}[Why orientation matters]
\label{rem:bf-ledh-ot-orientation}
If a derivation uses \(\Pi\) but the code applies \(T X\), transpose errors are
easy to introduce.  The VJP formulas below use \(T\), not \(\Pi\).  A reader who
wants the coupling form can substitute \(T=N\Pi^\top\).
\end{remark}

\section{VJP of the barycentric product}
\label{sec:bf-ledh-ot-barycentric-vjp}

The simplest and most important adjoint identity is the VJP of the barycentric
matrix product.

\begin{proposition}[Barycentric VJP]
\label{prop:bf-ledh-ot-barycentric-vjp}
Let \(\widetilde X=T X\), with \(T\in\mathbb R^{N\times N}\) and
\(X\in\mathbb R^{N\times d}\).  Use the Frobenius inner product for matrix
adjoints, and let \(G_{\widetilde X}\in\mathbb R^{N\times d}\) be the upstream
adjoint for \(\widetilde X\).  Then the direct adjoints are
\begin{equation}
\label{eq:bf-ledh-ot-vjp-barycentric}
  G_X^{\mathrm{direct}} = T^\top G_{\widetilde X},
  \qquad
  G_T = G_{\widetilde X}X^\top .
\end{equation}
\end{proposition}

\begin{proof}
Taking a differential gives
\[
  d\widetilde X = dT\,X + T\,dX .
\]
Pairing with \(G_{\widetilde X}\) under the Frobenius inner product,
\begin{align*}
  \langle G_{\widetilde X},d\widetilde X\rangle
  &=
  \langle G_{\widetilde X},dT\,X\rangle
  +
  \langle G_{\widetilde X},T\,dX\rangle  \\
  &=
  \langle G_{\widetilde X}X^\top,dT\rangle
  +
  \langle T^\top G_{\widetilde X},dX\rangle .
\end{align*}
The coefficients of \(dT\) and \(dX\) are exactly
\eqref{eq:bf-ledh-ot-vjp-barycentric}.
\end{proof}

Equation~\eqref{eq:bf-ledh-ot-vjp-barycentric} is only the direct part of the
particle adjoint.  Because \(T\) itself depends on \(X\) and \(\log w\), the
complete VJP also pulls \(G_T\) backward through the Sinkhorn or
annealed-transport construction.

\section{A normalized transport matrix}
\label{sec:bf-ledh-ot-normalized-transport}

For fixed final potentials \(f,g\in\mathbb R^N\), fixed normalized log weights
\(\log w_i\), and cost \(C_{ji}=c(z_j,z_i)\), the executed row-stochastic
transport can be written in column-normalized form
\begin{equation}
\label{eq:bf-ledh-ot-transport-form}
  T_{ji}
  =
  N w_i
  \frac{
    \exp\{(f_j+g_i-C_{ji})/\varepsilon\}
  }{
    \sum_{r=1}^N
    \exp\{(f_r+g_i-C_{ri})/\varepsilon\}
  } .
\end{equation}
This expression enforces \(T^\top\mathbf 1=Nw\) exactly for the displayed
normalization.  The row condition \(T\mathbf 1\approx\mathbf 1\) is enforced by
the potentials produced by finite Sinkhorn or annealed Sinkhorn iteration.
In the TensorFlow implementation this is the convention used by the final
transport application: \(f\) indexes output rows, \(g\) indexes source columns,
the normalization is a column log-sum-exp over output rows, and the transported
particles are obtained by the matrix product \(T X\).

Define
\begin{equation}
\label{eq:bf-ledh-ot-column-prob}
  p_{ji}
  =
  \frac{T_{ji}}{Nw_i},
  \qquad
  \sum_{j=1}^N p_{ji}=1.
\end{equation}
The vector \(p_{\cdot i}\) is the normalized column distribution associated with
source particle \(i\).

\begin{proposition}[VJP of the normalized transport output]
\label{prop:bf-ledh-ot-normalized-transport-vjp}
At fixed finite values of \(f,g\in\mathbb R^N\),
\(C\in\mathbb R^{N\times N}\), \(\varepsilon>0\), and normalized weights
\(w_i>0\), \(\sum_i w_i=1\), let \(T\) be defined by
\eqref{eq:bf-ledh-ot-transport-form}.  Let
\(G_T\in\mathbb R^{N\times N}\) be the upstream adjoint for \(T\).  For each
source column \(i\), set
\begin{equation}
\label{eq:bf-ledh-ot-column-mean-adjoint}
  \bar G_i=\sum_{r=1}^N p_{ri}(G_T)_{ri},
  \qquad
  S_{ji}=T_{ji}\bigl((G_T)_{ji}-\bar G_i\bigr).
\end{equation}
Then the VJP of \eqref{eq:bf-ledh-ot-transport-form} satisfies
\begin{align}
\label{eq:bf-ledh-ot-vjp-logw}
  (G_{\log w})_i
  &\mathrel{+}=
  \sum_{j=1}^N (G_T)_{ji}T_{ji},\\
\label{eq:bf-ledh-ot-vjp-f}
  (G_f)_j
  &\mathrel{+}=
  \frac{1}{\varepsilon}\sum_{i=1}^N S_{ji},\\
\label{eq:bf-ledh-ot-vjp-g}
  (G_g)_i
  &\mathrel{+}=
  \frac{1}{\varepsilon}\sum_{j=1}^N S_{ji},\\
\label{eq:bf-ledh-ot-vjp-cost}
  (G_C)_{ji}
  &\mathrel{+}=
  -\frac{1}{\varepsilon}S_{ji}.
\end{align}
If \(\log w\) is itself produced by a \(\operatorname{logsoftmax}\), then
\eqref{eq:bf-ledh-ot-vjp-logw} must be composed with the
\(\operatorname{logsoftmax}\) VJP.
\end{proposition}

\begin{proof}
Let
\[
  a_{ji}=\frac{f_j+g_i-C_{ji}}{\varepsilon},
  \qquad
  p_{ji}=\frac{\exp(a_{ji})}{\sum_r\exp(a_{ri})},
  \qquad
  T_{ji}=Nw_i p_{ji}.
\]
For a perturbation of one column \(i\),
\[
  dT_{ji}
  =
  T_{ji}
  \left[
    d\log w_i
    + da_{ji}
    - \sum_{r=1}^N p_{ri}da_{ri}
  \right].
\]
Pairing this expression with \((G_T)_{ji}\) and collecting the coefficient of
\(da_{ji}\) gives \(S_{ji}=T_{ji}((G_T)_{ji}-\bar G_i)\).  Since
\[
  da_{ji}
  =
  \frac{1}{\varepsilon}
  (df_j+dg_i-dC_{ji}),
\]
the adjoints for \(f\), \(g\), and \(C\) are
\eqref{eq:bf-ledh-ot-vjp-f}--\eqref{eq:bf-ledh-ot-vjp-cost}.  The coefficient
of \(d\log w_i\) is \(\sum_j (G_T)_{ji}T_{ji}\), giving
\eqref{eq:bf-ledh-ot-vjp-logw}.
\end{proof}

This proposition is the core OT VJP product.  It says that the upstream matrix
\(G_T\) should first be column-centered under the transport-induced probability
\(p_{\cdot i}\).  Only after that centering should the adjoint be pushed into
the potentials and costs.  In the displayed column-normalized map,
\(\sum_j S_{ji}=0\), so the final-potential adjoint \(G_g\) is identically zero
at this last normalization stage.  This is not a bug: adding the same \(g_i\)
to every logit in one normalized column does not change that column's
probabilities.  Earlier Sinkhorn-loop adjoints may still involve the variables
that produced \(g\).

\section{VJP of the quadratic cost}
\label{sec:bf-ledh-ot-cost-vjp}

For the common squared-distance cost
\[
  C_{ji}=\frac{1}{2}\|z_j-z_i\|^2,
\]
where the same scaled particle cloud \(Z\) supplies the row and column points,
the cost adjoint has another simple form.

\begin{proposition}[Quadratic-cost VJP]
\label{prop:bf-ledh-ot-cost-vjp}
Let \(G_C\) be the upstream adjoint for
\(C_{ji}=\frac12\|z_j-z_i\|^2\), where
\(Z=(z_1,\ldots,z_N)^\top\in\mathbb R^{N\times d}\).  The contribution to the
adjoint of \(z_k\in\mathbb R^d\) is
\begin{equation}
\label{eq:bf-ledh-ot-vjp-cost-points}
  (G_Z)_k
  \mathrel{+}=
  \sum_{i=1}^N (G_C)_{ki}(z_k-z_i)
  -
  \sum_{j=1}^N (G_C)_{jk}(z_j-z_k).
\end{equation}
\end{proposition}

\begin{proof}
The differential of one cost entry is
\[
  dC_{ji}=(z_j-z_i)^\top(dz_j-dz_i).
\]
Multiplying by \((G_C)_{ji}\), summing over \(j,i\), and collecting the
coefficient of \(dz_k\) gives the two sums in
\eqref{eq:bf-ledh-ot-vjp-cost-points}: the first from terms where \(k\) is the
row index, and the second from terms where \(k\) is the column index.
\end{proof}

If centering, scaling, or key construction is stopped by design, the VJP stops
there too.  If the selected mode differentiates through centering or scaling,
then \eqref{eq:bf-ledh-ot-vjp-cost-points} must be further composed with the
VJP of those preprocessing maps.

\section{Softmin and Sinkhorn adjoints}
\label{sec:bf-ledh-ot-softmin-vjp}

The finite Sinkhorn or annealed-transport step is a loop of softmin operations.
The exact schedule can differ between a fixed-target Sinkhorn route and the
FilterFlow-style annealed route, but the local adjoint primitive is the same.

\begin{definition}[Softmin block]
\label{def:bf-ledh-ot-softmin}
For a query index \(j\), costs \(C_{jr}\), values \(v_r\), and temperature
\(\varepsilon>0\), define
\begin{equation}
\label{eq:bf-ledh-ot-softmin}
  m_j
  =
  -\varepsilon
  \log
  \sum_{r=1}^N
  \exp\{(v_r-C_{jr})/\varepsilon\}.
\end{equation}
\end{definition}

\begin{proposition}[Softmin VJP for fixed temperature]
\label{prop:bf-ledh-ot-softmin-vjp}
Assume \(\varepsilon>0\) is held fixed and all displayed exponentials are
finite.  Let \(G_m\in\mathbb R^N\) be the upstream adjoint for \(m\), and define
\[
  q_{jr}
  =
  \frac{\exp\{(v_r-C_{jr})/\varepsilon\}}
       {\sum_s \exp\{(v_s-C_{js})/\varepsilon\}}.
\]
For fixed \(\varepsilon\),
\begin{align}
\label{eq:bf-ledh-ot-softmin-vjp-values}
  (G_v)_r
  &\mathrel{+}=
  -\sum_{j=1}^N (G_m)_j q_{jr},\\
\label{eq:bf-ledh-ot-softmin-vjp-costs}
  (G_C)_{jr}
  &\mathrel{+}=
  (G_m)_j q_{jr}.
\end{align}
\end{proposition}

\begin{proof}
Differentiating \eqref{eq:bf-ledh-ot-softmin} gives
\[
  dm_j
  =
  -\sum_r q_{jr}(dv_r-dC_{jr})
  =
  \sum_r q_{jr}dC_{jr}-\sum_r q_{jr}dv_r.
\]
Pairing with \(G_m\) and collecting coefficients yields
\eqref{eq:bf-ledh-ot-softmin-vjp-values} and
\eqref{eq:bf-ledh-ot-softmin-vjp-costs}.
\end{proof}

The sign and scale in Proposition~\ref{prop:bf-ledh-ot-softmin-vjp} are for a
potential-scale value \(v\).  Some implementations, including the TensorFlow
helper used in the current BayesFilter route, call the same operation with a
log-scale input \(h_r\):
\[
  m_j=-\varepsilon\log\sum_r
  \exp\{h_r-C_{jr}/\varepsilon\}.
\]
This is the same formula with \(v_r=\varepsilon h_r\).  Therefore the local VJP
with respect to \(h_r\) is \((G_h)_r\mathrel{+}=
-\varepsilon\sum_j(G_m)_j q_{jr}\), while the VJP with respect to a
potential-scale component \(v_r\) remains
\eqref{eq:bf-ledh-ot-softmin-vjp-values}.  If \(h_r=\log\alpha_r+b_r/\varepsilon\),
then the adjoint to \(b_r\) again has the potential-scale coefficient
\(-\sum_j(G_m)_j q_{jr}\).

For an annealed schedule, the loop state contains several potentials, for
example \((a_y,b_x,a_x,b_y)\).  Abstractly write the finite loop as
\begin{equation}
\label{eq:bf-ledh-ot-loop-state}
  s_{k+1}=F_k(s_k;Z,\log w,\varepsilon_k),
  \qquad k=0,\ldots,K-1.
\end{equation}
The reverse loop is then
\begin{equation}
\label{eq:bf-ledh-ot-reverse-loop}
  (G_{s_k},G_Z,G_{\log w})
  \mathrel{+}=
  D F_k(s_k;Z,\log w,\varepsilon_k)^\top G_{s_{k+1}},
  \qquad k=K-1,\ldots,0.
\end{equation}
Each multiplication by \(D F_k^\top\) is implemented by applying the softmin
VJP above, the averaging VJP for damped updates, and the cost VJP for the
pairwise squared distances.  This is ordinary reverse-mode differentiation of a
finite program, but expressed as an explicit adjoint scan rather than as a raw
tape over the whole program.

\begin{proposition}[Finite-Sinkhorn VJP by reverse scan]
\label{prop:bf-ledh-ot-finite-sinkhorn-vjp}
Assume the finite transport block is the deterministic program
\[
  (T,s_K)=\mathcal S_K(Z,\log w)
\]
obtained by a fixed number of differentiable softmin, averaging, cost, and
normalization operations with fixed branch decisions.  Then a custom VJP for
\(\mathcal S_K\) is obtained by:
\begin{enumerate}[label=(\roman*)]
  \item applying Proposition~\ref{prop:bf-ledh-ot-normalized-transport-vjp} to
  move \(G_T\) into final potentials, costs, and log weights;
  \item applying Proposition~\ref{prop:bf-ledh-ot-cost-vjp} to move final-cost
  adjoints into \(Z\);
  \item scanning \eqref{eq:bf-ledh-ot-reverse-loop} backward through the finite
  potential updates;
  \item composing any resulting \(G_Z\) with the chosen centering, scaling, and
  stop-gradient policy.
\end{enumerate}
The result equals the VJP of the executed finite program.
\end{proposition}

\begin{proof}
The finite transport block is a composition of differentiable maps under the
fixed branch and fixed iteration assumptions.  Propositions
\ref{prop:bf-ledh-ot-normalized-transport-vjp},
\ref{prop:bf-ledh-ot-cost-vjp}, and
\ref{prop:bf-ledh-ot-softmin-vjp} give the VJP for the nontrivial primitive
maps.  The VJP of a composition is obtained by applying the transposed local
Jacobians in reverse order.  That reverse composition is exactly the scan
described above.
\end{proof}

\section{Why a custom VJP is memory efficient}
\label{sec:bf-ledh-ot-memory}

The naive reverse graph for the OT block stores pairwise costs, logits,
normalizers, potentials, and intermediate loop states for every time step.  With
\(N\) particles and \(K\) Sinkhorn updates, this can scale like \(O(KN^2)\) in
stored tensors before the rest of the filtering graph is considered.

The custom VJP changes the storage problem.  The primal path must retain or
reconstruct only enough information to replay the finite program:
\begin{itemize}
  \item the input particles and log weights entering the OT block;
  \item the branch choices, masks, iteration count, temperature schedule, and
  stabilization policy;
  \item final potentials, and optionally sparse checkpoints of intermediate
  potentials;
  \item any stopped quantities needed to reproduce the executed scalar.
\end{itemize}
During the backward pass, pairwise blocks can be recomputed and reduced
immediately.  This trades additional arithmetic for lower peak memory.  The
trade is natural in HMC, because one scalar log target needs one reverse VJP to
obtain all parameter-score components, while a forward-mode JVP would need one
pass per requested direction.

\section{Embedding the OT VJP in the filtering loop}
\label{sec:bf-ledh-ot-filter-loop}

The full LEDH-PFPF-OT gradient is a reverse scan through time.  At time \(t\),
the adjoint entering the OT block contains two sources: downstream dependence
of future particles, and direct dependence of the current scalar contribution.
The OT VJP returns adjoints for the pre-transport particles and log weights:
\[
  (G_{X_t},G_{\log w_t})
  =
  \operatorname{VJP}_{\mathrm{OT},t}
  (G_{\widetilde X_t}).
\]
The log-weight adjoint then flows through the PF-PF proposal-correction terms:
transition density, observation density, proposal density, and flow
log-determinant.  The particle adjoint flows through the LEDH migration and the
model transition and observation callbacks.  Repeating this from \(t=T\) down
to \(t=1\) yields the parameter score.

\begin{proposition}[Filtering-loop custom VJP]
\label{prop:bf-ledh-ot-filtering-loop-vjp}
Suppose each time-step map in the executed LEDH-PFPF-OT value path is
differentiable under fixed branch decisions, fixed random streams, fixed masks,
and fixed iteration budgets.  If each time-step custom VJP is the VJP of that
same executed time-step map, then the reverse time scan returns the derivative
of the executed scalar \(\mathcal L_K(\theta)\).
\end{proposition}

\begin{proof}
The full filter value is a finite composition of the time-step maps and the
final scalar reduction.  By the chain rule, the VJP of the full composition is
the product of the time-step transposed Jacobians in reverse time order.  The
assumption says that each custom time-step VJP equals the corresponding
transposed Jacobian multiplication for the executed map.  Therefore their
reverse-time composition equals the VJP of the executed scalar.
\end{proof}

\section{Implementation contract}
\label{sec:bf-ledh-ot-implementation-contract}

A robust implementation should expose two separable functions:
\[
  \operatorname{ot\_forward}(X,\log w;\rho)
  =
  (\widetilde X,\mathrm{aux}),
  \qquad
  \operatorname{ot\_vjp}(\mathrm{aux},G_{\widetilde X})
  =
  (G_X,G_{\log w}),
\]
where \(\rho\) contains the declared OT settings: cost convention, epsilon,
annealing schedule, iteration budget, row and column chunk sizes, stabilization,
and stop-gradient policy.  The auxiliary record should contain only mathematical
or numerical data needed to reproduce the same finite map.  It should not
silently change the scalar by switching from streaming to dense semantics, from
finite Sinkhorn to another optimizer, or from one random stream to another.

The current TensorFlow implementation follows this contract through an
opt-in mode named \texttt{same\_scalar\_custom\_vjp}.  In dense mode, the
custom-gradient primal calls the same finite FilterFlow-style annealed
transport routine as the raw TensorFlow path, with the selected
\texttt{transport\_ad\_mode} and any declared warm-start state.  Its backward
rule then recomputes that same transport matrix under a local
\texttt{GradientTape} and applies the unclipped upstream transport adjoint.  In
streaming mode, the primal returns the transported particles without returning
a dense transport matrix; the backward rule recomputes the same streaming
transport under a local \texttt{GradientTape}.  This is a same-scalar
checkpointed or recomputed VJP route.  It is not yet a fully hand-coded
reverse scan of every softmin and normalization primitive in
Proposition~\ref{prop:bf-ledh-ot-finite-sinkhorn-vjp}.  The hand-derived
reverse scan remains the mathematical destination for a future lower-level
adjoint implementation.

The following checks belong with any implementation:
\begin{enumerate}[label=(\arabic*)]
  \item row and column residuals for the finite transport object;
  \item equality of dense and streaming VJP on small problems where both are
  feasible;
  \item finite-difference or raw-autodiff agreement on small fixed-seed cases;
  \item explicit tests for stopped versus differentiated scale, keys, and
  potentials;
  \item static-shape behavior under the intended compiled path;
  \item finite value and finite score on the target model class before any HMC
  use.
\end{enumerate}

These checks certify local derivative accounting for the declared finite
computation.  They do not certify the statistical accuracy of the relaxed
filter, posterior convergence, or default production readiness.

\section{Relation to forward-mode JVP}
\label{sec:bf-ledh-ot-jvp}

Forward-mode JVP is also mathematically legitimate.  For a perturbation
\(\dot\theta\), it propagates tangents through the same filtering loop and
returns
\[
  D_\theta\mathcal L_K(\theta)\,\dot\theta .
\]
This can be memory efficient, and it is useful for diagnostics along a small
number of directions.  For HMC, however, the sampler usually needs the whole
score vector.  If the parameter dimension is \(d_\theta\), forward mode needs
roughly \(d_\theta\) directional passes unless it is batched over tangent
directions.  A VJP gives the full scalar-target score in one reverse pass.  That
is why the custom-gradient design in this chapter emphasizes the OT VJP.

\section{What this chapter does and does not claim}
\label{sec:bf-ledh-ot-boundary}

This chapter gives a mathematical route for a custom gradient of the executed
finite LEDH-PFPF-OT scalar.  It claims:
\begin{itemize}
  \item the barycentric OT product has the VJP in
  \eqref{eq:bf-ledh-ot-vjp-barycentric};
  \item the normalized finite-transport matrix has the VJP in
  \eqref{eq:bf-ledh-ot-vjp-logw}--\eqref{eq:bf-ledh-ot-vjp-cost};
  \item finite Sinkhorn or annealed Sinkhorn can be differentiated by a reverse
  scan through the executed softmin program;
  \item composing those local VJPs through the filtering loop gives the
  derivative of the same executed scalar under fixed branch and randomness
  assumptions.
\end{itemize}

It does \emph{not} claim that the relaxed OT filter equals a categorical
particle filter, that finite Sinkhorn equals unregularized OT, that a custom VJP
improves Monte Carlo accuracy, or that the resulting score is automatically
HMC-ready for a structural model.  Those are separate validation questions.

\FloatBarrier
